% ============================================
% Example: Methods from Hu et al., Nature Ecology & Evolution 2025
% "Collective dynamical regimes predict invasion success and impacts in microbial communities"
% DOI: 10.1038/s41559-024-02618-y
% ============================================
% NOTE: This file is for reference only - to guide the style and structure
% of our Methods section for Nature Ecology & Evolution submission.
% ============================================

% Key observations from Hu et al. 2025:
%
% 1. METHODS SECTION STRUCTURE:
%    - Methods appear at END of main text (after Results/Discussion)
%    - Referenced as "Online Methods" in the paper
%    - Detailed protocols go to Supplementary Information
%
% 2. EXTENDED DATA:
%    - 10 Extended Data Figures used for essential background
%    - Extended Data Fig. 1: Taxonomic identity of isolates (phylogeny)
%    - Extended Data Fig. 2-7: Invasion outcome matrices and time series
%    - Extended Data Fig. 8-10: Relative species abundances
%
% 3. SUPPLEMENTARY INFORMATION:
%    - Contains: Materials and Methods, Notes, Discussion, Figs. 1-33
%    - Full experimental details in "Supplementary Materials and Methods"
%    - 33 Supplementary Figures for additional analyses
%
% 4. MAIN TEXT METHODS - what's included:
%    - Brief reference to ASV identification via 16S sequencing
%    - Essential experimental parameters only
%    - Equations for theoretical models (likely in main text or SI)
%    - References to Extended Data and Supplementary figures
%
% 5. DATA/CODE AVAILABILITY (separate sections at end):
%    - Data availability: "Isolates and communities are available upon request.
%      All data are available in the Supplementary Information and via Dryad"
%    - Code availability: "All codes used for simulation and analysis in this
%      publication are available via GitHub"
%
% ============================================
% IMPLICATIONS FOR OUR MANUSCRIPT:
% ============================================
%
% A. METHODS should be CONCISE - main experimental approach only
%    - NO full buffer recipes
%    - NO detailed parameter lists
%    - Reference SI for "see Supplementary Methods for details"
%
% B. EQUATIONS can stay in Methods OR move to SI
%    - Core model equation (gLV) can stay in main Methods
%    - Classification equations should move to SI
%
% C. Use EXTENDED DATA for:
%    - Phylogenetic tree (currently Suppl. Fig. S1)
%    - Time series data
%    - Outcome matrices
%
% D. SUPPLEMENTARY INFORMATION structure:
%    - Supplementary Materials and Methods (detailed protocols)
%    - Supplementary Notes (extended analysis)
%    - Supplementary Discussion (if needed)
%    - Supplementary Figures (numbered S1, S2, etc.)
%
% ============================================
